\documentclass[12pt]{article}
\usepackage[utf8]{inputenc}
\usepackage[T1]{fontenc}
\usepackage{hyperref}
\title{Kayu Agung Language Resources for K-12 Curriculum Development}
\author{Compiled by Claude Sonnet 5 (claude-sonnet-5) for Andr\'as Kornai}
\date{}
\begin{document}
\maketitle

\section{Introduction}
Kayu Agung (vky, Glottocode kayu1244) is a retired ISO 639-3 code, merged into Komering (kge) in 2008. Glottolog files it under ``Bookkeeping'' and notes high intelligibility with upstream Komering varieties, and two speech varieties in Kayu Agung town. Speaker numbers in sources conflict and are not given here.

\section{Primary Educational Resources}
None specific to Kayu Agung were found.

\section{Linguistic and Reference Materials}
\begin{itemize}
\item Glottolog entry: \url{https://glottolog.org/resource/languoid/id/kayu1244}
\item Verbix page on Komering: \url{https://docs.verbix.com/Languages/Komering}
\item Conference paper on South Sumatra languages: \url{https://proceeding.unnes.ac.id/snpasca/article/download/2260/1743/5944}
\end{itemize}

\section{Religious and Community Resources}
Ethnographic profiles: \url{https://joshuaproject.net/people_groups/12597} and \url{https://www.prayingforindonesia.com/ethnic-groups/the-people-of-sumatra/who-are-the-kayu-agung/}

\section{Digital and Audio Resources}
None found.

\section{Curriculum Development Context}
Because the code is retired and the variety treated as Komering, curriculum work should probably target Komering and check with local speakers how Kayu Agung varieties relate to it.
\end{document}